\section{Introduction}
\label{sec:intro}

Most reinforcement-learning environments that contain a hazard put the hazard
somewhere the learner can see it directly. Monitoring and suppression tasks
put it \emph{in the reward}: coverage of a fire front, mapping error,
suppression progress. Safe-RL formulations put it \emph{in a constraint}: a
cost channel, a Lagrangian penalty, a safety critic. In both lineages the
hazard is an object the agent is asked to reason about. This paper studies
the remaining case, which we argue is the common one for robots in damaged
structures: the hazard is neither target nor constraint but \emph{weather}. It
evolves under its own dynamics, it kills, and nothing in the objective names
it.

The Compound Hostile Environment (\CHE) makes that case precise. A
cellular-automaton fire spreads over a grid while a team of agents performs a
task whose reward reads task variables only. The fire is coupled to the rest
of the world in two directions. Structural collapse events seed new ignitions
(Coupling A), so the environment manufactures its own danger; and the smoke
field attenuates what agents can observe by Beer--Lambert transmittance
(Coupling B), so perception fails exactly where and when the hazard is
active. A third axis, communication denial, degrades the inter-agent link
graph independently of the hazard. From a single agent's point of view the
hazard looks like nonstationarity. Proposition~\ref{prop:markov} says it is
not: it is state under a stationary kernel, and every standard multi-agent
training tool applies unmodified.

\paragraph{The compound question.}
Stressors that co-occur in the world interact causally: collapse ignites
fire, fire makes smoke, smoke blinds. Does a policy need to \emph{see} them
co-occur during training to survive their co-occurrence at test time? The
question is not new; it has priority holders in the compositional-generalization
literature (\S\ref{sec:related}). What has been missing is an instrument clean
enough to answer it. The obvious comparison---one policy trained on the
stressors jointly against one trained on them separately---confounds two
things: whether the policy saw the stressors \emph{together}, and how
\emph{often} it saw each. In our own confirmatory contrast the jointly trained
arm sees each element three times as often as the isolated arm. We therefore
run that contrast as registered, and beside it a \emph{matched-marginal dose
sweep} in which each element's frequency is held at exactly one half while
the co-occurrence fraction moves from zero to one half. The sweep is the
comparison that standard joint-versus-isolated training studies do not run,
and it is what distinguishes this work from the multi-task and
domain-randomization literature that asks joint-versus-isolated questions
routinely. In the vocabulary of distribution-based compositionality
assessment~\citep{keysers2020}, the sweep drives compound divergence while
holding atom divergence at zero; the endpoint contrast does not, and we say
so.

\paragraph{Why an instrument first.}
A single pre-registered experiment has four registered outcomes, and three of
them are not the founding hypothesis. We designed the paper so that a paper
exists on every branch: the environment, its calibration, and a measurement
methodology for empirical RL are the contributions that do not depend on the
sign of one contrast. Stating this before the result is what makes the
pre-registration credible rather than decorative.

\paragraph{Contributions.}
\begin{enumerate}[leftmargin=*,itemsep=2pt]
  \item \textbf{The environment as an instrument} (\S\ref{sec:env}). A
  factored Dec-POMDP whose reward is provably hazard-independent
  (Definition~\ref{def:ambient}), whose step order is a theorem's hypothesis
  (Proposition~\ref{prop:markov}), whose ablations are bitwise-exact nested
  models, and whose severity levels are defined by a measured dynamical phase
  rather than by knob settings.
  \item \textbf{Four calibration findings} (\S\ref{sec:severity}--%
  \S\ref{sec:comms}) that stand without any training comparison: the
  collapse-to-fire channel self-limits where the fire is worst; compound
  hostility is rare and bursty; perception decay concentrates on danger
  moments and cannot be behaviourally suppressed; and communication is inert
  under a frozen encoder at swarm scale.
  \item \textbf{A measurement methodology} (\S\ref{sec:method}): per-metric,
  per-hardware, per-artifact reproducibility floors; contrasts graded on the
  contrast's own standard error; design-stage power stated as an 80\%-power
  minimum detectable effect at the family-corrected level; and the
  distinction between recording a channel and grading it.
  \item \textbf{A pre-registered, mechanically blinded test of compositional
  generalization} (\S\ref{sec:test}), with $k = 72$ seeds per arm on a single
  card, a matched-marginal dose sweep, and its result. \gap{one clause from
  the branch file}.
\end{enumerate}

\paragraph{What we do not claim.}
No convergence guarantee for the training loop; no theorem that joint
training wins (Proposition~\ref{prop:shift} bounds what can be guaranteed,
not what happens); and no exact critical value for the implemented kernel,
whose phase structure is measured rather than assumed. The confirmatory
estimand is a contrast at \emph{matched budget}, never at convergence.
